\section{引言：语言模型能推理吗？}

本节课分为两大主题：（1）语言模型的推理能力——能否用语言模型进行数学、逻辑、空间等领域的多步推理；（2）语言模型智能体（Agents）——能否让语言模型在真实环境中执行操作、完成任务。这两个主题紧密相关：智能体需要在环境中进行多步推理和规划才能完成复杂任务。

\begin{figure}[H]
\centering
\includegraphics[width=0.95\textwidth]{slides-images/slide-000.jpg}
\caption{CS224N Lecture 14 封面：Reasoning and Agents，讲者 Shikhar Murty\protect\footnotemark}
\end{figure}
\footnotetext{Slides 第1页。}

\begin{knowledgebox}{课程背景}
本讲内容基于近3--4年的最新研究，许多问题尚无定论。讲者特别提醒：关于``语言模型是否真的在推理''这一核心问题，学术界仍存在大量争论，本讲将展示正反两方面的证据。在前几讲（Lecture 9--11）中，我们已经了解到 LLM 擅长生成符合人类偏好的文本延续，本讲的核心问题是：\textbf{它们是否也能进行推理？}
\end{knowledgebox}

\subsection{什么是推理}

推理（Reasoning）的核心定义是：\textbf{使用事实（facts）和逻辑（logic）来得出结论}。

\begin{figure}[H]
\centering
\includegraphics[width=0.95\textwidth]{slides-images/slide-004.jpg}
\caption{推理的定义与演绎推理示例：从前提``所有哺乳动物都有肾脏''和``所有鲸鱼都是哺乳动物''推出``所有鲸鱼都有肾脏''\protect\footnotemark}
\end{figure}
\footnotetext{Slides 第5页。Slide credit: Graham Neubig (11-711 ANLP)。}

推理可以分为三种基本类型：

\begin{importantbox}{推理的三种类型}
\begin{itemize}
    \item \textbf{演绎推理（Deductive Reasoning）}：从逻辑规则和前提出发，推导出\textbf{确定性}结论。例如：所有哺乳动物有肾脏 $+$ 所有鲸鱼是哺乳动物 $\Rightarrow$ 所有鲸鱼有肾脏。
    \item \textbf{归纳推理（Inductive Reasoning）}：从观察到的实例中总结规律，得出\textbf{可能性}结论。例如：每次看到有翅膀的生物都是鸟 $\Rightarrow$ 看到有翅膀的生物，推测它可能是鸟。
    \item \textbf{溯因推理（Abductive Reasoning）}：从观察到的现象出发，推测可能的\textbf{解释}。例如：车无法启动且引擎下有液体 $\Rightarrow$ 可能是散热器漏液。
\end{itemize}
\end{importantbox}

\subsection{形式推理与非形式推理}

\begin{figure}[H]
\centering
\includegraphics[width=0.95\textwidth]{slides-images/slide-007.jpg}
\caption{形式推理 vs 非形式推理：形式推理遵循公理和逻辑规则；非形式推理依赖直觉、经验和常识\protect\footnotemark}
\end{figure}
\footnotetext{Slides 第8页。}

\begin{itemize}
    \item \textbf{形式推理（Formal Reasoning）}：遵循形式逻辑的公理和规则来推导真值条件。
    \item \textbf{非形式推理（Informal Reasoning）}：利用直觉、经验和常识来得出结论，这是我们日常生活中最常用的推理方式。
\end{itemize}

本讲中提到的``推理''主要指\textbf{非形式的演绎推理}，且通常涉及\textbf{多个步骤}。

\subsection{本章小结}

推理是利用事实和逻辑得出结论的过程，分为演绎、归纳和溯因三类。本讲聚焦于非形式多步推理，探讨语言模型能否通过提示或训练来展现推理能力，以及这种能力的本质和局限性。

%% ========== Part 2: 提示方法 ==========

